% Additive Lane 4 proof repairs.  This file is not included by a canonical
% manuscript and carries no theorem promotion by itself.

\section{Dependency and descent repairs for the quartic audit}

\subsection{The plane degree criterion and its field of application}

The plane result used after a coordinate has been straightened is the theorem
of Appelgate and Onishi:

\begin{quote}
J. Appelgate and H. Onishi, ``The Jacobian conjecture in two variables,''
\emph{Journal of Pure and Applied Algebra} \textbf{37} (1985), 215--227,
DOI \texttt{10.1016/0022-4049(85)90099-4}.  The criterion is stated in the
paper's opening theorem/abstract formulation: a complex plane Keller pair is
invertible when the degree of at least one coordinate has at most two prime
factors, counted with multiplicity.
\end{quote}

The lemmas used in that argument must not be cited as if their original proofs
were unproblematic.  The repair chain is:

\begin{quote}
A. Nowicki and Y. Nakai, ``On Appelgate--Onishi's Lemmas,''
\emph{Journal of Pure and Applied Algebra} \textbf{51} (1988), no.~3,
305--310, DOI \texttt{10.1016/0022-4049(88)90069-2}.  Lemmas A and B are
stated on p.~305, the characteristic-zero base field is fixed in \S1 on
p.~306, and the corrected reformulations/proofs are in \S4, pp.~309--310.
See also A. Nowicki and Y. Nakai, ``Correction to `On Appelgate--Onishi's
lemmas','' \emph{Journal of Pure and Applied Algebra} \textbf{58} (1989),
no.~1, 101.
\end{quote}

Here is the field bookkeeping needed in Lane~4.  Put $K=k(t)$ and let
$\overline K$ be an algebraic closure.  After straightening the known
coordinate, the remaining pair $f,g\in\overline K[u,v]$ has constant nonzero
Jacobian and one coordinate degree in
\[
\{1,2,3,4,5,6,7,9\}.
\]
Each integer in this set has at most two prime factors with multiplicity.
The Appelgate--Onishi argument, with the Nowicki--Nakai repairs, is algebraic
in characteristic zero; equivalently one may spread the finitely many
coefficients to a finitely generated subfield of characteristic zero, embed
that field in $\mathbb C$, apply the complex theorem, and descend the finite
polynomial identities.  Thus the criterion applies over $\overline K$.

Let $(U,V)\in\overline K[s_1,s_2]^2$ be the inverse.  The inverse is unique.
For every $\sigma\in\operatorname{Aut}(\overline K/K)$, the pair
$(\sigma U,\sigma V)$ is another inverse to the same $K$-defined pair, hence
$\sigma U=U$ and $\sigma V=V$.  Therefore $U,V\in K[s_1,s_2]$.  This is the
field descent actually used by the three-variable reduction.  A later appeal
to the birational Keller implication is a separate imported theorem and must
not be folded into the plane degree criterion.

\subsection{Expanded homogeneous cubic centralizer argument}

\begin{lemma}[Homogeneous cubic centralizer with descent]
Let $k$ be a characteristic-zero field, let $x,y,z$ be independent variables,
and let $R\in k[x,y,z]$ be homogeneous of degree three.  Assume
\[
x\nmid R
\quad\text{and}\quad
R\notin k[x,\ell]
\quad\text{for every nonzero }\ell\in k[y,z]_1.
\]
Set $K=k(x)$ and
\[
\delta=J_{y,z}(R,-):K[y,z]\longrightarrow K[y,z].
\]
Then
\[
\ker\delta=K[R].
\]
Moreover, if $Q\in k[x,y,z]$ is homogeneous of degree $d$ and
$\delta(Q)=0$, then
\[
Q=\sum_{0\le j\le d/3}c_jx^{d-3j}R^j,
\qquad c_j\in k.
\]
In particular the polynomial kernel on homogeneous elements is generated by
$x$ and $R$.
\end{lemma}

\begin{proof}
First regard $R$ as a polynomial in $K[y,z]$.  A nontrivial compositional
decomposition $R=u(S)$ with $\deg u>1$ would, because the total degree in
$y,z$ is at most three, force the homogeneous leading part of $R$ to be a
power of a $K$-linear form.  Clearing denominators and comparing the
homogeneous $x$-degrees then puts $R$ in $k[x,\ell]$ for a nonzero
$\ell\in k[y,z]_1$, contrary to the hypothesis.  Thus $R$ is a closed
polynomial over $K$.  The two-variable Jacobian centralizer/common-generator
theorem in characteristic zero gives
\[
\ker J_{y,z}(R,-)=K[R].
\]

Let $Q$ be homogeneous of degree $d$ in this kernel.  There is a unique finite
expansion
\[
Q=\sum_j f_j(x)R^j,\qquad f_j(x)\in k(x),
\]
because $R$ is transcendental over $K$.  Apply the scalar action
$(x,y,z)\mapsto(\lambda x,\lambda y,\lambda z)$.  Since $Q$ has degree $d$
and $R$ has degree three, uniqueness of the $K[R]$ expansion yields
\[
f_j(\lambda x)=\lambda^{d-3j}f_j(x)
\quad\text{in }k(\lambda,x).
\]
A one-variable rational function satisfying this identity is a monomial:
putting $x=1$ gives $f_j(\lambda)=f_j(1)\lambda^{d-3j}$, hence
\[
f_j(x)=c_jx^{d-3j},\qquad c_j\in k.
\]

It remains to exclude negative exponents.  The assumption $x\nmid R$ says
that $R$, considered in $k[x][y,z]$, is primitive with respect to the prime
$x$.  Gauss' lemma (or the $x$-adic valuation) gives the intersection
identity
\[
k(x)[R]\cap k[x,y,z]=k[x,R].
\]
Indeed, if a coefficient in the unique $R$-expansion had a pole of maximal
order along $x=0$, multiplying by that power of $x$ and reducing modulo $x$
would give a nonzero polynomial relation among powers of
$R(0,y,z)$ over $k$, impossible because $R(0,y,z)$ is nonconstant.  Since
$Q$ belongs to the intersection, every $f_j$ is in $k[x]$.  Therefore
$d-3j\ge0$ whenever $c_j\ne0$, and the displayed formula follows.
\end{proof}

This expanded proof separates the three inputs that were compressed in the
short packet argument: closedness/common generation over $k(x)$, homogeneity
of one-variable rational coefficients, and polynomial descent across the
prime $x=0$.

\subsection{Pointed normal form for the degree-five fixed-factor basepoint}

The basepoint condition below belongs to the degree-five/six appendix, not to
the quartic case tree.

\begin{proposition}[Residual normal form for the quintic $(g,e)=(1,2)$ boundary]
Let
\[
H_5=G(A^2,AB,B^2)
\]
with $G$ linear and $A,B$ homogeneous quadrics.  Assume that $A,B$ form a
regular sequence, that $A/B$ is nonconstant on the line $G=0$, and that
$V_{\mathbf P^2}(G,A,B)$ is nonempty.  After a projective source change, a
constant change of the pair $(A,B)$, and choosing one basepoint, one may write
\[
G=z,\qquad p=[1:0:0],
\]
\[
A=xy+z\ell_1,
\qquad
B=y^2+z\ell_2,
\qquad \ell_1,\ell_2\in k[x,y,z]_1.
\]
The remaining coefficient space must still satisfy $\gcd(A,B)=1$ and the
Keller jet equations.
\end{proposition}

\begin{proof}
Send the line $G=0$ to $z=0$ and the chosen basepoint to $[1:0:0]$.  The
restrictions $a=A|_{z=0}$ and $b=B|_{z=0}$ are binary quadrics vanishing at
that point, so both are divisible by $y$.  They cannot both be multiples of
$y^2$, because then their ratio on $z=0$ would be constant wherever defined.
After replacing $(A,B)$ by an invertible constant linear combination, their
residual linear factors are independent.  A projective change preserving the
point and the line makes those residual factors $x$ and $y$.  Hence
$a=xy$ and $b=y^2$ up to nonzero scalars, which are absorbed into the target
change.  Restoring the terms divisible by $z$ gives the displayed form.

This is a normalization of the surviving basepoint locus, not an exclusion.
The basepoint-free Hensel/Koszul lifting theorem cannot be applied because
$(A,B)$ is not a regular sequence in the homogeneous coordinate ring of
$G=0$ at $p$.
\end{proof}
